\subsection{Etap Plan i Code}
Pierwsza faza cyklu DevOps obejmuje planowanie wymagań oraz wytwarzanie kodu źródłowego. W niniejszym projekcie etapy te są ze sobą ściśle zintegrowane, co pozwala na pełną identyfikowalność zmian — od zgłoszenia zadania po konkretną linię kodu.

\subsubsection{Zarządzanie projektami: Jira}
Do zarządzania wymaganiami i planowania prac wykorzystano platformę \textbf{Jira} (projekt \textit{INZ}). Narzędzie to wybrano ze względu na jego status standardu rynkowego oraz zaawansowane możliwości wsparcia metodyk zwinnych (Agile).

\textbf{Metodyka:}
Agile to rodzina iteracyjnych metodyk wytwarzania oprogramowania, kładących nacisk na częste dostarczanie wartości i adaptację do zmian. Sprinty są podstawową jednostką iteracji w Agile — to ograniczone czasowo okresy pracy (zazwyczaj 1--4 tygodnie), po których dostarcza się działający przyrost produktu. Tablica Kanban to uproszczona alternatywa skupiona na wizualizacji przepływu zadań i ograniczaniu pracy w toku, bez sztywnych iteracji czasowych.

Ponieważ projekt realizowany jest indywidualnie i w ograniczonym czasie, zadania podzielono logicznie na sprinty zgodnie z etapami pętli DevOps, lecz do bieżącego śledzenia pracy zastosowano tablicę Kanban — bardziej odpowiednią dla małych, jednoosobowych projektów.

\textbf{Struktura backlogu:}
Backlog projektu odwzorowuje etapy pętli DevOps. Zdefiniowano pięć Epics odpowiadających sprintom implementacyjnym: \textit{Build}, \textit{Release}, \textit{Deploy}, \textit{Monitor} i \textit{Polish}. W ramach każdego Epica utworzono zestaw zadań (\textit{Tasks}) pokrywających zarówno dostarczenie kodu infrastruktury, jak i uzupełnienie odpowiednich sekcji pracy inżynierskiej. Łącznie backlog liczy 29 zadań.

\textbf{Integracja i dobre praktyki:}
Zastosowano integrację z systemem kontroli wersji GitHub. Każda gałąź kodu (\textit{branch}) jest nazwana według schematu \texttt{feature/INZ-XX-opis}, a wiadomości zatwierdzeń zawierają identyfikator zadania (np. \texttt{feat(helm): INZ-9 create umbrella chart}). Dzięki temu zmiany w kodzie są automatycznie powiązane z odpowiednim zgłoszeniem w Jira, co zapewnia pełną identyfikowalność — od wymagania po konkretną linię kodu.

\begin{figure}[H]
    \centering
    \includegraphics[width=0.75\textwidth]{jira.png}
    \caption{Tablica Kanban projektu INZ w Jira -- widok sprintów i zadań}
    \label{fig:jira-board}
\end{figure}

\subsubsection{System kontroli wersji: GitHub}
Jako system kontroli wersji i platformę współpracy wybrano \textbf{GitHub}, oparty na rozproszonym systemie Git. Jest on centralnym punktem prawdy (\textit{Single Source of Truth}) dla całego monorepozytorium.

\textbf{Standard Conventional Commits:}
Zgodnie z decyzją projektową \textit{ADR-0004}, wprowadzono standard \textbf{Conventional Commits}. Polega on na nadawaniu wiadomościom zatwierdzeń ustrukturyzowanej formy (np. \texttt{feat:}, \texttt{fix:}, \texttt{chore:}).
\begin{itemize}
    \item \textbf{Uzasadnienie:} Standaryzacja pozwala na automatyczne generowanie list zmian (\textit{changelogs}) oraz automatyzację wersjonowania semantycznego (\textit{Semantic Versioning}).
    \item \textbf{Rola w CI/CD:} Czytelna historia zmian umożliwia narzędziom automatyzującym precyzyjne określenie, które komponenty systemu wymagają przebudowania i wdrożenia po każdej zmianie w monorepozytorium.
\end{itemize}

Poniższy przykład ilustruje kompletną wiadomość zatwierdzenia zgodną ze standardem. Nagłówek zawiera typ, zakres oraz zwięzły opis zmiany, a stopka (\textit{footer}) powiązuje commit z konkretnym zadaniem w systemie Jira:

\begin{lstlisting}[caption={Przykład wiadomości zatwierdzenia zgodnej ze standardem Conventional Commits}, label={lst:conventional-commit}]
feat(helm): create umbrella chart for all 7 services

Aggregates all microservice charts under a single parent
chart, enabling one-command deployment of the full stack.

Resolves: INZ-9
\end{lstlisting}

\subsubsection{Wytwarzanie oprogramowania: Stos technologiczny}
Proces kodowania realizowany jest przy użyciu nowoczesnych środowisk programistycznych (\textit{IntelliJ IDEA}, \textit{VS Code}), a sama implementacja opiera się na technologiach wspierających paradygmat mikroserwisów.

\textbf{Backend (Java i Quarkus — tryb GraalVM Native Image):}
Mikroserwisy backendowe zbudowano z użyciem frameworka \textit{Quarkus} skompilowanego do natywnych plików wykonywalnych za pomocą \textit{GraalVM Native Image} (kompilacja \textit{Ahead-of-Time}, AOT). W trybie natywnym Quarkus analizuje ścieżki wykonania kodu w czasie kompilacji i wyklucza nieużywane klasy, co skutkuje:
\begin{itemize}
    \item \textbf{Błyskawicznym startem} — aplikacja jest gotowa do obsługi żądań w czasie poniżej 100 ms, co skraca czas wznowienia po awarii i przyspiesza skalowanie poziome w Kubernetes.
    \item \textbf{Niskim zużyciem pamięci} — brak środowiska uruchomieniowego JVM zmniejsza footprint kontenera, co bezpośrednio przekłada się na efektywność kosztową w środowiskach chmurowych.
    \item \textbf{Minimalnym obrazem kontenera} — obraz oparty na \texttt{ubi-minimal} zawiera wyłącznie system operacyjny i skompilowany binarny, bez JDK ani bibliotek środowiska uruchomieniowego.
\end{itemize}
Wstrzykiwanie zależności realizowane jest przez mechanizm CDI Quarkusa (\textit{ArC}), który w trybie AOT przetwarza wszystkie adnotacje w czasie kompilacji, a nie w czasie uruchomienia.

\textbf{Frontend (Vue 3 + Vite):}
Zgodnie z decyzją \textit{ADR-0014}, warstwę prezentacji zbudowano z użyciem frameworka \textit{Vue 3} i narzędzia budującego \textit{Vite}. Kod pisany jest w \textit{TypeScript}, a poprawność typów weryfikuje \textit{vue-tsc} przed każdym buildem produkcyjnym. Komponenty interfejsu komunikują się z mikroserwisami backendowymi przez REST API.

\textbf{Pryncypia Twelve-Factor App:}
Wszystkie komponenty systemu są projektowane zgodnie z zasadami \textit{Twelve-Factor App}. Kluczowe aspekty to:
\begin{itemize}
    \item \textbf{Konfiguracja:} Cała konfiguracja specyficzna dla środowiska jest wstrzykiwana przez zmienne środowiskowe, a nie zaszyta w kodzie.
    \item \textbf{Procesy:} Aplikacje są bezstanowe (\textit{stateless}), co umożliwia ich łatwe skalowanie w poziomie wewnątrz klastra Kubernetes.
    \item \textbf{Logowanie:} Logi są traktowane jako strumień zdarzeń wysyłany na standardowe wyjście, co umożliwia ich centralizację przez narzędzia observability.
\end{itemize}
